% Extracted missing sections from Chapter 3
% Auto-generated from OCR cleaned files

\section[Introduction to the Scholarship]{\origsecnum{3.1}Introduction to the Scholarship}
\label{sec:chapters-ch03:3.1}

Deciphering
the Inscriptions
3.1
Introduction<!-- FOOTNOTE:1 -->For the famous story, first published by Tung in 1930, about Wang Yi-jung's sickness and Liu E's discovery of inscriptions on the ``dragon bones'' from a Peking pharmacy, see Lefeuvre (1975), p. 17, n. 1; Li Chi (1977), pp. 7-11. Liu E himself made no mention of the incident in his introduction to T'ieh-yün.<!-- /FOOTNOTE --> to the Scholarship
Few events have excited the minds of modern Chinese scholars like the discovery
of the oracle bones. In the three-quarters of a century that has elapsed since 1899, when
Wang Yi-jung and Liu E first appreciated the true nature of these ``dragon bone'' inscriptions, a large body of books and articles-at the last count, over 180 books and over 970
monographs has appeared “like bamboo shoots after spring rain.'' In this chapter, I
will first introduce the essential works and indicate how they may be used.<!-- FOOTNOTE:2 -->Shima (1971), preface, p. 1; cf. Chang Pingch'üan (1967), p. 841, who refers to approximately 1,000 articles and over 180 books by about 300 scholars.<!-- /FOOTNOTE -->³<!-- FOOTNOTE:3 -->The progress of oracle-bone scholarship (chiaku-hsüeh) has involved several interrelated stages: the identification of the true nature of the inscriptions; the identification of the site from which the inscribed bone fragments were being taken; the collection and publication of those fragments; the eventual conduct of scientific excavations in the Hsiao-t'un region; and, above all, the gradual decipherment, still continuing, of the graphs. The history of this progress has been well summarized elsewhere. See, in particular, Kaizuka (1946), pp. 191-201; (1967), pp. 8-60; Jung Keng (1947); Hu Hou-hsüan (1955), pp. 35-42; Ch'en Meng-chia (1956), pp. 50-54; Tung (1965), pp. 46-118. For Western-language accounts, see Creel (1938), pp. 1-16; Lefeuvre (1975), which is fundamental; and Li Chi (1977).3.2
Published Sources
Approximately 107,000 oracle-bone pieces have now been assembled in over
eighty collections, generally in the form of rubbings.<!-- FOOTNOTE:4 -->Matsumaru (1973), pp. 19-22, lists eightyeight collections (excluding those like Ping-pien, Cho-ts'un, and T'ung-tsuan, which duplicate material previously published) and indicates the style 57 ] U L U<!-- /FOOTNOTE -->4 A list of the major collections is

\section[Reference Works]{\origsecnum{3.3}Reference Works}
\label{sec:chapters-ch03:3.3}

Reference Works
Over<!-- FOOTNOTE:7 -->Kaizuka (1967, p. 204, estimates that we know about 3,000 graphs on the oracle bones. about 800 of which can be interpreted; Chang Ping-ch'üan (1967), p. 829, refers to 3,320 graphs, of which we can recognize about 1,000, with 500 to 700 more identifiable but in dispute. Hu Houhsüan (1945), p. 5b, speculated, extravagantly in my view, that the full oracle-bone vocabulary, were we able to recover it, would have included 5,000 to 6,000 words, of which we would be able to recognize from 2,000 to over 3,000. No exactitude is possible in such counts (for other examples, see Chang Ch'un-shu [1972], p. 283); nor is anything more than a general estimate necessary. See too n. 20.<!-- /FOOTNOTE --> three thousand Shang graphs have been identified on the oracle bones, and
about half of these may be translated or identified with varying degrees of certainty. The
dictionaries and concordances listed below are indispensable to the study of the inscriptions.
3.3.1
Dictionaries<!-- FOOTNOTE:8 -->Chung-kuo k'e-hsüeh-yüan k'ao-ku yen-chiu-so ‹† 國科學院考古研究所,Chia-ku wen-pien 甲骨文編, K'ao-ku-hsüeh chuan-k'an yi-chung ti-shih-ssu hao 考古學專刋乙種第十四號(Peking,1965). This work, prepared under the general editorship of Sun Hai-po is a revised edition of his Chia-ku wen-pien (Peiping, 1934), which identified some 500 to 600 graphs, taken from eight collections (cf. Teng and Biggerstaff [1971], p. 46). Some of the 4,672 entries in CKWP probably refer to the same words and could be combined. The 1,723 entries in the main section of the dictionary consist of words which have been ``identified'' either in the sense that a Shuo-wen graph can be assigned (in 941 cases) or that a modern graphic reconstruction and putative position within the Shuo-wen system of classification can be provided. A ho-wen section contains 371 entries for graphs which frequently appear in conjunction, such as those used in ancestral titles. The appendix contains 2,949 entries of graphs which have not been deciphered or classified (CKWP, preface, p. 1). The forty collections consulted are listed with their abbreviations between the appendix and the modern-character index. An edition of this dictionary was reissued by Yi-wen in Taipei in 1975 under the title Chiao-cheng chia-ku wen-pien 校正甲骨文編,8
The standard dictionary is Chia-ku wen-pien (abbreviated ChIP), which contains
a total of 4,672 entries taken from forty collections of published inscriptions. This needs
to be supplemented, on occasion, with Hsü chia-ku wen-pien.''
The entries in both dictionaries, like those in many works of oracle-bone scholarship,
are arranged according to the 540 radicals of Shuo-wen. Each entry consists of the Shuo-wen
graph (when known) as heading, followed by hand-copied reproductions of the oracle-bone characters it is thought to represent. Occasionally, there is brief discussion of the
graph's meaning, but generally the Shuo-wen graph is the only “definition” provided. The

[^9]: Chin Hsiang-heng, Hsü chia-ku wenpien (Taipei, 1959). This is a supplement to Sun Hai-po's 1934 dictionary (referred to in n. 8). Based on thirty-eight published collections, it contains 1,048 oracle-bone graphs identified with Shuo-wen equivalents, 1,585 graphs not found in Shuo-wen, an appendix of ho-wen, and an appendix of some 300 unidentified graphs. There are two modern-character indexes for the graphs with and without Shuo-wen equivalents. Cf. Noel Barnard in RBS 5 (1959), p. 463. All three dictionaries may need to be consulted on occasion. For example, for the graph yüanë, CKWP (1965) has more citations (thirty-two) than the 1934 edition (nine citations) and the Hsü chia-ku wen-pien (twelve citations) combined. But due to the decision of the 1965 editors (preface, p. 3) not to reproduce identical forms of the same character, the 1965 edition has many less entries for the graph chiu than the 1934 edition. 59 11 } U 11 0

\subsection[Dictionaries]{\origsecnum{3.3.1}Dictionaries}
\label{sec:chapters-ch03:3.3.1}

Reference Works
Over<!-- FOOTNOTE:7 -->Kaizuka (1967, p. 204, estimates that we know about 3,000 graphs on the oracle bones. about 800 of which can be interpreted; Chang Ping-ch'üan (1967), p. 829, refers to 3,320 graphs, of which we can recognize about 1,000, with 500 to 700 more identifiable but in dispute. Hu Houhsüan (1945), p. 5b, speculated, extravagantly in my view, that the full oracle-bone vocabulary, were we able to recover it, would have included 5,000 to 6,000 words, of which we would be able to recognize from 2,000 to over 3,000. No exactitude is possible in such counts (for other examples, see Chang Ch'un-shu [1972], p. 283); nor is anything more than a general estimate necessary. See too n. 20.<!-- /FOOTNOTE --> three thousand Shang graphs have been identified on the oracle bones, and
about half of these may be translated or identified with varying degrees of certainty. The
dictionaries and concordances listed below are indispensable to the study of the inscriptions.
3.3.1
Dictionaries<!-- FOOTNOTE:8 -->Chung-kuo k'e-hsüeh-yüan k'ao-ku yen-chiu-so ‹† 國科學院考古研究所,Chia-ku wen-pien 甲骨文編, K'ao-ku-hsüeh chuan-k'an yi-chung ti-shih-ssu hao 考古學專刋乙種第十四號(Peking,1965). This work, prepared under the general editorship of Sun Hai-po is a revised edition of his Chia-ku wen-pien (Peiping, 1934), which identified some 500 to 600 graphs, taken from eight collections (cf. Teng and Biggerstaff [1971], p. 46). Some of the 4,672 entries in CKWP probably refer to the same words and could be combined. The 1,723 entries in the main section of the dictionary consist of words which have been ``identified'' either in the sense that a Shuo-wen graph can be assigned (in 941 cases) or that a modern graphic reconstruction and putative position within the Shuo-wen system of classification can be provided. A ho-wen section contains 371 entries for graphs which frequently appear in conjunction, such as those used in ancestral titles. The appendix contains 2,949 entries of graphs which have not been deciphered or classified (CKWP, preface, p. 1). The forty collections consulted are listed with their abbreviations between the appendix and the modern-character index. An edition of this dictionary was reissued by Yi-wen in Taipei in 1975 under the title Chiao-cheng chia-ku wen-pien 校正甲骨文編,8
The standard dictionary is Chia-ku wen-pien (abbreviated ChIP), which contains
a total of 4,672 entries taken from forty collections of published inscriptions. This needs
to be supplemented, on occasion, with Hsü chia-ku wen-pien.''
The entries in both dictionaries, like those in many works of oracle-bone scholarship,
are arranged according to the 540 radicals of Shuo-wen. Each entry consists of the Shuo-wen
graph (when known) as heading, followed by hand-copied reproductions of the oracle-bone characters it is thought to represent. Occasionally, there is brief discussion of the
graph's meaning, but generally the Shuo-wen graph is the only “definition” provided. The

[^9]: Chin Hsiang-heng, Hsü chia-ku wenpien (Taipei, 1959). This is a supplement to Sun Hai-po's 1934 dictionary (referred to in n. 8). Based on thirty-eight published collections, it contains 1,048 oracle-bone graphs identified with Shuo-wen equivalents, 1,585 graphs not found in Shuo-wen, an appendix of ho-wen, and an appendix of some 300 unidentified graphs. There are two modern-character indexes for the graphs with and without Shuo-wen equivalents. Cf. Noel Barnard in RBS 5 (1959), p. 463. All three dictionaries may need to be consulted on occasion. For example, for the graph yüanë, CKWP (1965) has more citations (thirty-two) than the 1934 edition (nine citations) and the Hsü chia-ku wen-pien (twelve citations) combined. But due to the decision of the 1965 editors (preface, p. 3) not to reproduce identical forms of the same character, the 1965 edition has many less entries for the graph chiu than the 1934 edition. 59 11 } U 11 0

\subsection[Concordances and Compendiums]{\origsecnum{3.3.2}Concordances and Compendiums}
\label{sec:chapters-ch03:3.3.2}

dictionary. 14 Following the order of the graphs in Chin-wen-pien (1959), it not only attempts
to reproduce every occurrence of the graph in question that may be found in the published
corpus of Shang and Chou vessels,<!-- FOOTNOTE:14 -->Chou Fa-kao with Chang Jih-sheng [Cheung Yat-shing], Hsü Chih-yi [Tsui Chee-yee] and Lin Chich-ming [Lam Kit-ming]林潔明,eds., Chin-wen ku-lin 金文詰林, 16 vols. (Hong Kong, 1974-1975). For a review, see Mattos (1976).<!-- /FOOTNOTE --><!-- FOOTNOTE:15 -->Jung Keng's original corpus in Chin-wenpien (1959) has been expanded to include fiftythree more inscriptions from Kuo Mo-jo (1957).<!-- /FOOTNOTE --><!-- FOOTNOTE:16 -->Bernhard Karlgren, Grammata Serica Recensa (Stockholm, 1957).<!-- /FOOTNOTE --><!-- FOOTNOTE:17 -->Tōdō Akiyasu, Kanji gogen jiten (Tokyo, 1969). For a review of an earlier version of this work, see Pulleyblank (1968).<!-- /FOOTNOTE --><!-- FOOTNOTE:18 -->Katō Joken, Kanji no kigen 起原(Tokyo, 1970).<!-- /FOOTNOTE --><!-- FOOTNOTE:19 -->Shima Kunio, Inkyo bokuji sōrui , 2d. rev. ed. (Tokyo, 1971). The first edition (Tokyo, 1967), which is now less useful because it lacks the modern-character index, modern-character identifications at the entry heads, in-text cross references, cross references to CKIT, and Yi-pien/Ping-pien identifications of the second edition, was reprinted in Taipei by T'aishun shu-chuin 1970. Despite the reduced page size of the second edition (1971), it is to be preferred in every respect to the first and is the edition usually cited in this book.<!-- /FOOTNOTE --><!-- FOOTNOTE:20 -->From the table of contents (S9-12) I estimate that the work contains about 3,095 entries. Shima, however, indicates that the 840 entries for which cross references are given to CKWT constitute no more than one-fifth of the entries (front matter, p. 3), so he would seem to place the number of entries at over 4,000. 61 ㄈㄈ U } 0 ☐ ☐<!-- /FOOTNOTE --> 15 it also quotes the sentence in which the graph appears.
Some twenty thousand of these inscription sentences are reproduced in the 1,894 entries.
In addition, commentaries by other scholars dealing with the graph or the sentence are
quoted in full, together with copious notations by Chou Fa-kao and his fellow editors.
Reproducing the inscriptional context and assembling the scholarship in this way, Chin-wen ku-lin is a major reference work for etymological study.
The fact that these five dictionaries are all arranged by the Shuo-wen system has
advantages. In the first place, the 540 radicals provide a convenient series of headings
under which to place graphs for which no modern form exists. Further, since the dictionaries are organized on the same principle, a scholar may move rapidly from one to
another; a graph at the start of chian 5 in one will be at the start of chüan 5 in all (it is for
this reason that old style citations are preferable to new [for example, CKWP, ch. 12,
p. 18b, rather than CKWP, p. 496]). The modern-character indexes at the rear permit
ready entry into the system for those not familiar with it.
Since the study of the inscriptions is not just the study of graphic forms but is also
the study of language (cf. sec. 3.5.2 on phonology), it goes without saying that a student
of the Shang inscriptions will consult other etymological dictionaries as the need arises.
Particularly useful are Grammata Serica Recensa (abbreviated GSR), 16 Kanji gogen jiten, 17
and Kanji no kigen. 18 These dictionaries attempt to establish phonetic and semantic word
families in early Chinese, and sometimes refer to the oracle inscriptions.
Concordances and Compendiums
Inkyo bokuji sōrui (abbreviated S) is a concordance of the oracle inscriptions in
sixty-three collections. 19 Over three thousand words have been classified. 20 Under each

\section[How to Read the Graphs on a Bone or Shell Fragment]{\origsecnum{3.6}How to Read the Graphs on a Bone or Shell Fragment}
\label{sec:chapters-ch03:3.6}

How to Read the Graphs on a
Bone or Shell Fragment
Nothing appears more arcane, at first glance, than an oracle-bone rubbing. I
will try to demystify it in this section with a step-by-step analysis of how an actual turtle-shell piece, composed of two fragments (fig. 20), may be identified, deciphered, and placed
in context. The flow chart (fig. 21: summarizes the possible approaches.
How to Identify the Fragment
Orienting the piece is rarely a problem since rubbings are normally printed
correctly.<!-- FOOTNOTE:64 -->Fragments have been printed upside down, especially in the earliest collections. See, for example, T'ieh-yün 23.4, 28.4, and the other cases listed in Yen Yi-p'ing's preface to the Yi-wen (Taipei, 1959) reprint of T'ieh-yün. Errors of this sort appear with some frequency in secondary works, e.g., Eisenberger (1938), p. 67 (the inscription is probably fake); Ch'en Chih-mai (1966), fig. 4a (whole plastron inverted); William S. Y. Wang (1972), p. 53 (bottom two fragments, left column, inverted).<!-- /FOOTNOTE --> 64 If necessary, orientation is best done either by identifying certain graphs, or
graph elements, whose correct orientation is known, or by orienting the crack numbers
(near the top of the pu crack; sec. 2.4) and crack notations (near the bottom of the pu
crack; sec. 2.6) correctly. In figure 20, the crack number 6 and the shang-chi at the
left are readily placed. The pu cracks all run left; and the crack numbers 6, 7, and 8 run
right; these clues, together with the general configuration of the piece, lead us to suppose
that it comes from the right side of a shell in the region of the hyoplastron or hypoplastron
(fig.<!-- FOOTNOTE:65 -->See ch. 2, n. 122.<!-- /FOOTNOTE --> 3).65
As we have seen, inscriptions were written from top to bottom (sec. 2.9.4), but it is
necessary to determine in this case, as in all others, whether the inscription should be read
sideways and then down (in fig. 20, graphs 1, 4, 6, 2, 3, 5, 7; or 6, 4, 1, 2, 7, 5, 3) or read
down, with the columns read left to right (graphs 6, 7, 4, 5, 1, 2, 3) or right to left (graphs
I to 7 read in normal sequence).<!-- FOOTNOTE:66 -->There is also the possibility that the graphs on a fragment do not belong to the same inscription (as defined at ch. 2, n. 1) but are parts of two or more separate inscriptions; as we shall see, that is not the case here.<!-- /FOOTNOTE -->66 The fact that inscriptions generally run against the
cracks (sec. 2.9.4) suggests that the inscription in this case should be read from left to right,
but, as we see almost immediately from the position of the preface, we are dealing with a
common exception.<!-- FOOTNOTE:67 -->See ch. 2, n. 126. 0 71 U CO<!-- /FOOTNOTE -->67
The first, virtually automatic, step in identifying and translating any inscription is
basically for negating intransitive/passive verbs
and *pjat for negating transitive/causative verbs.
He also finds that the *P-type negative was used
when the situation involved natural forces or other
phenomena beyond Shang control, while the
*M-type was used when an element of human
will or control was involved. Nivison (1977) provides support for some of these conclusions. Mickel
(1976) has undertaken a similar study of the
``disaster'' words of the Wu Ting period (cf. ch.
2, n. 80, above).
64. Fragments have been printed upside down,
especially in the earliest collections. See, for
example, T'ieh-yün 23.4, 28.4, and the other cases
listed in Yen Yi-p'ing's preface to the Yi-wen
(Taipei, 1959) reprint of T'ieh-yün. Errors of this
sort appear with some frequency in secondary
works, e.g., Eisenberger (1938), p. 67 (the inscription is probably fake); Ch'en Chih-mai
(1966), fig. 4a (whole plastron inverted); William
S. Y. Wang (1972), p. 53 (bottom two fragments,
left column, inverted).
65. See ch. 2, n. 122.
66. There is also the possibility that the graphs
on a fragment do not belong to the same inscription
(as defined at ch. 2, n. 1) but are parts of two or
more separate inscriptions; as we shall see, that
is not the case here.
67. See ch. 2, n. 126.

\subsection[How to Identify the Fragment]{\origsecnum{3.6.1}How to Identify the Fragment}
\label{sec:chapters-ch03:3.6.1}

How to Read the Graphs on a
Bone or Shell Fragment
Nothing appears more arcane, at first glance, than an oracle-bone rubbing. I
will try to demystify it in this section with a step-by-step analysis of how an actual turtle-shell piece, composed of two fragments (fig. 20), may be identified, deciphered, and placed
in context. The flow chart (fig. 21: summarizes the possible approaches.
How to Identify the Fragment
Orienting the piece is rarely a problem since rubbings are normally printed
correctly.<!-- FOOTNOTE:64 -->Fragments have been printed upside down, especially in the earliest collections. See, for example, T'ieh-yün 23.4, 28.4, and the other cases listed in Yen Yi-p'ing's preface to the Yi-wen (Taipei, 1959) reprint of T'ieh-yün. Errors of this sort appear with some frequency in secondary works, e.g., Eisenberger (1938), p. 67 (the inscription is probably fake); Ch'en Chih-mai (1966), fig. 4a (whole plastron inverted); William S. Y. Wang (1972), p. 53 (bottom two fragments, left column, inverted).<!-- /FOOTNOTE --> 64 If necessary, orientation is best done either by identifying certain graphs, or
graph elements, whose correct orientation is known, or by orienting the crack numbers
(near the top of the pu crack; sec. 2.4) and crack notations (near the bottom of the pu
crack; sec. 2.6) correctly. In figure 20, the crack number 6 and the shang-chi at the
left are readily placed. The pu cracks all run left; and the crack numbers 6, 7, and 8 run
right; these clues, together with the general configuration of the piece, lead us to suppose
that it comes from the right side of a shell in the region of the hyoplastron or hypoplastron
(fig.<!-- FOOTNOTE:65 -->See ch. 2, n. 122.<!-- /FOOTNOTE --> 3).65
As we have seen, inscriptions were written from top to bottom (sec. 2.9.4), but it is
necessary to determine in this case, as in all others, whether the inscription should be read
sideways and then down (in fig. 20, graphs 1, 4, 6, 2, 3, 5, 7; or 6, 4, 1, 2, 7, 5, 3) or read
down, with the columns read left to right (graphs 6, 7, 4, 5, 1, 2, 3) or right to left (graphs
I to 7 read in normal sequence).<!-- FOOTNOTE:66 -->There is also the possibility that the graphs on a fragment do not belong to the same inscription (as defined at ch. 2, n. 1) but are parts of two or more separate inscriptions; as we shall see, that is not the case here.<!-- /FOOTNOTE -->66 The fact that inscriptions generally run against the
cracks (sec. 2.9.4) suggests that the inscription in this case should be read from left to right,
but, as we see almost immediately from the position of the preface, we are dealing with a
common exception.<!-- FOOTNOTE:67 -->See ch. 2, n. 126. 0 71 U CO<!-- /FOOTNOTE -->67
The first, virtually automatic, step in identifying and translating any inscription is
basically for negating intransitive/passive verbs
and *pjat for negating transitive/causative verbs.
He also finds that the *P-type negative was used
when the situation involved natural forces or other
phenomena beyond Shang control, while the
*M-type was used when an element of human
will or control was involved. Nivison (1977) provides support for some of these conclusions. Mickel
(1976) has undertaken a similar study of the
``disaster'' words of the Wu Ting period (cf. ch.
2, n. 80, above).
64. Fragments have been printed upside down,
especially in the earliest collections. See, for
example, T'ieh-yün 23.4, 28.4, and the other cases
listed in Yen Yi-p'ing's preface to the Yi-wen
(Taipei, 1959) reprint of T'ieh-yün. Errors of this
sort appear with some frequency in secondary
works, e.g., Eisenberger (1938), p. 67 (the inscription is probably fake); Ch'en Chih-mai
(1966), fig. 4a (whole plastron inverted); William
S. Y. Wang (1972), p. 53 (bottom two fragments,
left column, inverted).
65. See ch. 2, n. 122.
66. There is also the possibility that the graphs
on a fragment do not belong to the same inscription
(as defined at ch. 2, n. 1) but are parts of two or
more separate inscriptions; as we shall see, that
is not the case here.
67. See ch. 2, n. 126.

\section[Reading Practices]{\origsecnum{3.7}Reading Practices}
\label{sec:chapters-ch03:3.7}

There is no sure Ariadne to guide us. CKWT, for instance, is not comprehensive;
Jimbun's glosses are not complete and are out of date. No gloss is infallible. Scholars have
to plunge into the labyrinth for themselves, using the chart provided by figure 21. The
best guide despite occasional errors, variant transcriptions, and omissions-is Shima's
Sōrui. Sōrui assembles the inscriptions; it guides us to CKWT and to the modern glosses;
it shows how a word was used in context. The rest we must do for ourselves. And we do it
by developing philological expertise, by studying the examples of graphic and semantic
evolution that commentators have proposed, and, above all, by reading as widely as
possible in the inscriptions themselves.
A Plastron Set
Translation not only involves placing inscriptions in the context of other inscriptions. It also involves understanding the role that divination and the carving of the inscriptions played in Shang culture; it involves understanding, as far as we can, what the diviner
thought he was doing.<!-- FOOTNOTE:82 -->For a preliminary analysis of these issues, which will be considered more fully in Studies, see Keightley (1973a) and (1975c). Cf. too, Takashima (1976), who proposes a rationale for some of the charges in the set discussed below.<!-- /FOOTNOTE --> 82 It is instructive in this regard to examine one set of inscriptions
in detail in an initial attempt to explore the rationale underlying Shang divination and to
indicate the kinds of questions that we need to ask.
I select the set of inscriptions carved on the five plastrons Ping-pien 12-21 because,
when the charges on the front are considered together with the subcharges (sec. 3.7.4.2)
and prognostications on the back, they form one of the most informative multiple-plastron
sets we now possess. (The fact that these five plastrons form a set [sec. 2.5] is demonstrated
by the similarity of the divination charges and the ordinal sequence of crack numbers
[sec. 2.4], which are all 1 on Ping-pien 12, all 2 on Ping-pien 14, and so on.)
The analysis which follows should, if possible, be studied with the Ping-pien rubbings
and Chang Ping-chüan's k'ao-shih at one's side; I do not always follow his interpretation,
and I make no attempt to reproduce his extensive notes about the identity of the various
figures and statelets involved. My concerns are primarily with the divination process.
3.7.1
3.7.1.1
The Inscriptions
Duplication and Abbreviation
With the exception of minor variations in layout and epigraphy, inscriptions I,
2, 5, 8, 9, and 10 on the front of the shells and 1 through 8 on the back were ``carbon
copies,'' reproduced without abbreviation on all five plastrons.<!-- FOOTNOTE:83 -->This conclusion depends, in several cases, on supplying graphs that would, it is assumed, have appeared on missing shell fragments. Due to the roughness of the inner surface of the shell, it is not always easy to see the graphs in the rubbings of the plastron backs; in these cases I rely on the k'ao-shih.<!-- /FOOTNOTE -->83 The other inscriptions

\subsection[The Inscriptions]{\origsecnum{3.7.1}The Inscriptions}
\label{sec:chapters-ch03:3.7.1}

There is no sure Ariadne to guide us. CKWT, for instance, is not comprehensive;
Jimbun's glosses are not complete and are out of date. No gloss is infallible. Scholars have
to plunge into the labyrinth for themselves, using the chart provided by figure 21. The
best guide despite occasional errors, variant transcriptions, and omissions-is Shima's
Sōrui. Sōrui assembles the inscriptions; it guides us to CKWT and to the modern glosses;
it shows how a word was used in context. The rest we must do for ourselves. And we do it
by developing philological expertise, by studying the examples of graphic and semantic
evolution that commentators have proposed, and, above all, by reading as widely as
possible in the inscriptions themselves.
A Plastron Set
Translation not only involves placing inscriptions in the context of other inscriptions. It also involves understanding the role that divination and the carving of the inscriptions played in Shang culture; it involves understanding, as far as we can, what the diviner
thought he was doing.<!-- FOOTNOTE:82 -->For a preliminary analysis of these issues, which will be considered more fully in Studies, see Keightley (1973a) and (1975c). Cf. too, Takashima (1976), who proposes a rationale for some of the charges in the set discussed below.<!-- /FOOTNOTE --> 82 It is instructive in this regard to examine one set of inscriptions
in detail in an initial attempt to explore the rationale underlying Shang divination and to
indicate the kinds of questions that we need to ask.
I select the set of inscriptions carved on the five plastrons Ping-pien 12-21 because,
when the charges on the front are considered together with the subcharges (sec. 3.7.4.2)
and prognostications on the back, they form one of the most informative multiple-plastron
sets we now possess. (The fact that these five plastrons form a set [sec. 2.5] is demonstrated
by the similarity of the divination charges and the ordinal sequence of crack numbers
[sec. 2.4], which are all 1 on Ping-pien 12, all 2 on Ping-pien 14, and so on.)
The analysis which follows should, if possible, be studied with the Ping-pien rubbings
and Chang Ping-chüan's k'ao-shih at one's side; I do not always follow his interpretation,
and I make no attempt to reproduce his extensive notes about the identity of the various
figures and statelets involved. My concerns are primarily with the divination process.
3.7.1
3.7.1.1
The Inscriptions
Duplication and Abbreviation
With the exception of minor variations in layout and epigraphy, inscriptions I,
2, 5, 8, 9, and 10 on the front of the shells and 1 through 8 on the back were ``carbon
copies,'' reproduced without abbreviation on all five plastrons.<!-- FOOTNOTE:83 -->This conclusion depends, in several cases, on supplying graphs that would, it is assumed, have appeared on missing shell fragments. Due to the roughness of the inner surface of the shell, it is not always easy to see the graphs in the rubbings of the plastron backs; in these cases I rely on the k'ao-shih.<!-- /FOOTNOTE -->83 The other inscriptions

\subsubsection[Duplication and Abbreviation]{\origsecnum{3.7.1.1}Duplication and Abbreviation}
\label{sec:chapters-ch03:3.7.1.1}

There is no sure Ariadne to guide us. CKWT, for instance, is not comprehensive;
Jimbun's glosses are not complete and are out of date. No gloss is infallible. Scholars have
to plunge into the labyrinth for themselves, using the chart provided by figure 21. The
best guide despite occasional errors, variant transcriptions, and omissions-is Shima's
Sōrui. Sōrui assembles the inscriptions; it guides us to CKWT and to the modern glosses;
it shows how a word was used in context. The rest we must do for ourselves. And we do it
by developing philological expertise, by studying the examples of graphic and semantic
evolution that commentators have proposed, and, above all, by reading as widely as
possible in the inscriptions themselves.
A Plastron Set
Translation not only involves placing inscriptions in the context of other inscriptions. It also involves understanding the role that divination and the carving of the inscriptions played in Shang culture; it involves understanding, as far as we can, what the diviner
thought he was doing.<!-- FOOTNOTE:82 -->For a preliminary analysis of these issues, which will be considered more fully in Studies, see Keightley (1973a) and (1975c). Cf. too, Takashima (1976), who proposes a rationale for some of the charges in the set discussed below.<!-- /FOOTNOTE --> 82 It is instructive in this regard to examine one set of inscriptions
in detail in an initial attempt to explore the rationale underlying Shang divination and to
indicate the kinds of questions that we need to ask.
I select the set of inscriptions carved on the five plastrons Ping-pien 12-21 because,
when the charges on the front are considered together with the subcharges (sec. 3.7.4.2)
and prognostications on the back, they form one of the most informative multiple-plastron
sets we now possess. (The fact that these five plastrons form a set [sec. 2.5] is demonstrated
by the similarity of the divination charges and the ordinal sequence of crack numbers
[sec. 2.4], which are all 1 on Ping-pien 12, all 2 on Ping-pien 14, and so on.)
The analysis which follows should, if possible, be studied with the Ping-pien rubbings
and Chang Ping-chüan's k'ao-shih at one's side; I do not always follow his interpretation,
and I make no attempt to reproduce his extensive notes about the identity of the various
figures and statelets involved. My concerns are primarily with the divination process.
3.7.1
3.7.1.1
The Inscriptions
Duplication and Abbreviation
With the exception of minor variations in layout and epigraphy, inscriptions I,
2, 5, 8, 9, and 10 on the front of the shells and 1 through 8 on the back were ``carbon
copies,'' reproduced without abbreviation on all five plastrons.<!-- FOOTNOTE:83 -->This conclusion depends, in several cases, on supplying graphs that would, it is assumed, have appeared on missing shell fragments. Due to the roughness of the inner surface of the shell, it is not always easy to see the graphs in the rubbings of the plastron backs; in these cases I rely on the k'ao-shih.<!-- /FOOTNOTE -->83 The other inscriptions

\subsection[Translation]{\origsecnum{3.7.2}Translation}
\label{sec:chapters-ch03:3.7.2}

Translation
of Ping-pien 12, 14, 16, 18, 20
(Preface:)
(Charge:)
(Preface:)
(Charge:)
PAIRS OF COMPLEMENTARY CHARGES
Negative charge (left side)
(2) Crack-making on hsin-yu (day
58), Ch üeh divined:
``This season, the king should not
follow Wang Ch'eng to attack the Hsia
Wei. for if he does, we) will not perhaps receive assistance in this case.<!-- FOOTNOTE:85 -->Chang Ping-ch'üan reads the bone graph as ch'un #, ``spring.'' Chao Feng (1976), p. 222, reverting to an earlier theory of Yang Shu-ta , argues that the graph (composed, in his view, of k'ou □ plus ts'ai *) was a loan for tsai ``year.'' CKWT (pp. 1972-1974) rejects both readings and argues, with Yü Hsing-wu (1940), pp. 6a-8b, that the graph was ancestral to *d'iôg or *t'iôg t'iao, which is close in sound to *ts'iôg/ch'iu K, ``autumn.'' A study of all the inscriptions containing S187.2-4) reveals that the period to which it applied extended from at least the eleventh to the fifth Shang month. ``Season'' is offered as a functional translation. For the reason why I translate these charges in the obligatory mood, see n. 44. The phrase shou yu yu has been much discussed. There is little doubt that the final graph stands for yu, thiori, ``(divine) assistance.'' As to the 4 (or, in later oracle texts, * ; see table 26, no. 4.3), I follow the view that it has an attributive, pronominal function. This is the solution briefly suggested by scholars such as Wu Ch'i-ch'ang ([1971], pp. 10It) and Ch’en Meng-chia ([1956], pp. 96-97); it is well and extensively supported by David Nivison (1977). In this view, the shou yu yu of this inscription would mean ``receive this assistance'' (i.e., assistance in this attack against the Hsia Wei); shou yu nien would mean ``receive this harvest'' (i.e., the one divined about, as in fig. 9). For the alternative theory (discussed by Nivison [1977], pp. 18-19) that 4 in such phrases stands for yu #, ``abundant,'' see Takashima (1973), pp. 336-338; cf. Serruys (1974), pp. 29-30. For additional scholarship, see Chin Hsiang-heng (1967); Nivison (1973)%; Yen Yi-p'ing (1973a).<!-- /FOOTNOTE --><!-- FOOTNOTE:86 -->I assume, on the basis of context, that sentences 3 and 4 are in the same conditional mood as I and 2, though recorded in abbreviated form. This is speculative. If the meaning in parentheses is not supplied, then the sentences should read: ``The king should follow Chih Kuo/The king should not follow Chih Kuo.'' For ``should not'' as a translation for the negative wu, see n. 44. To supply the Pa-fang as the object of the attack is still more speculative. I do so for the following reasons. No inscription (S65.4 ff.) records that Chih Kuo was to attack the Hsia Wei. It is unlikely, therefore, that sentences 3 to 6 should be taken as proposing that the king join Chih Kuo in attacking the Hsia Wei. Ping-pien 22.11-12, probably divined by Ch'üch on keng-shen (day 57), concerns the king following Chih Kuo to attack the Pa-[fang]. Ping-pien 24 + Ching-chin 1266 (see Ping-pien 24, k'ao-shih, pp. 48-51) records eight charges (A-H), probably all divined by Cheng on hsin-yu (day 58), about the king following Wang Ch'eng to attack the Hsia Wei and twelve charges (I-N, P-U),<!-- /FOOTNOTE -->"85
(4) Crack-making on hsin-yu (day
58), Ch'üeh divined:
``(This season) the king should not
follow Chih Kuo (to attack the Pa-fang. for if he does, we will not perhaps
receive assistance in this case)."86
Positive charge right side)
(1) Crack-making on hsin-yu (day
58), Ch'üch divined:
``This season, the king should follow
Wang Ch'eng to attack the Hsia Wei,
(for if he does, we will receive assistance in this case.''
(3) Crack-making on hsin-yu (day
58), Ch'üeh divined:
``(This season) the king should follow
Chih Kuo (to attack the Pa-fang, for
if he does, we will receive assistance in
this case).''
Ping-ch'uan (1956), pp. 246, 253-254; Chou
Hung-hsiang (1969), pp. 60-69; Li Ta-liang
(1972), pp. 134-169. Li notes (pp. 137, 162-163)
that it is the later inscriptions in a set that are
likely to be abbreviated; but that is not true in
this case. He also notes pp. 138, 157) that when
three or more inscriptions refer to the same topic,
the preface and charge are both frequently
abbreviated. It is probable that in extreme cases,
the entire inscription could be abbreviated out of
existence; see ch. 2, n. 32, above, on inscriptionless
cracks.

\subsection[The Crack Notations]{\origsecnum{3.7.3}The Crack Notations}
\label{sec:chapters-ch03:3.7.3}

The Hollows and Cracks
Supplementing our ignorance about the piece missing from the left epiplastron
of Ping-pien 17 with information derived from the corresponding areas of Ping-pien 15, 19,
and 21, we can infer that each plastron had thirty-two hollows chiseled into its back, with
sixteen on the right and sixteen on the left, arranged in a symmetrical pattern (fig. 23).
A total of 160 hollows was made, therefore, on the five plastrons. A study of the pu cracks
appearing on the front of the plastrons suggests that the hollows were all burned and
cracked; the hollows were used with full efficiency (cf. sec.<!-- FOOTNOTE:92 -->For ``if'' as a translation of ch'i in prognostications, see n. 41.<!-- /FOOTNOTE --><!-- FOOTNOTE:93 -->The hsin is illegible in the rubbing. Here I follow the transcriptions of Chang Ping-ch'uan and S94.2.<!-- /FOOTNOTE --><!-- FOOTNOTE:94 -->Following Ping-pien 13, k'ao-shih, p. 36, Chih (or Sui 7 [?]-see CKWT, pp. 18951896) is a name that appears in four uninformative period I inscriptions (S61.2). Fei (?) is a hapax legomenon. Possibly, Chih Fei referred to a leader called Fei from the region of Chih. I assume with Shima (S357.4) that is a variant of . Chang Ping-ch'üan, in fact, transcribes the graph as ; this is possible on the basis of Ping-pien 13, but Shima's reading is supported by Ping-pien 21. The roughness of the plastron back makes the accuracy of the rubbings uncertain. As a matter of convenience, I follow Ikeda (1964), 2.15.14 and 2.32.14, for the highly speculative reading of Tuan (?). On this place, see n. 98 below.<!-- /FOOTNOTE --><!-- FOOTNOTE:95 -->The rubbings do not allow certainty on this point for they do not always reproduce the pu cracks, especially when the cracks were not engraved (cf. Chang Ping-ch'üan [1956], p. 241). Note, for example, that Ping-pien 20.3 has a crack number 5, but no crack is visible in the rubbing; similarly, the hsiao-chi crack notation on the left hypoplastron of Ping-pien 18 lacks a visible crack. The fact that cracks can be identified in each of the thirty-two hollow positions for at least one of the plastrons suggests that all the hollows were cracked in each case. This analysis, and much of what follows, depends on the use of tracings made of the hollow positions, reversed and superimposed on the front of the shell, to locate the area where cracks should have formed and near which the relevant inscriptions should have been carved. more<!-- /FOOTNOTE --><!-- FOOTNOTE:96 -->For example, two cracks without numbers appear at the top center of the right and left hypoplastrons of Ping-pien 16, and two appear at the top of the right and left xiphiplastrons; the cracks associated with Ping-pien 18.7 and 18.8 have no numbers (the k'ao-shih appears to be wrong), and at least two more inscriptionless cracks on the right hyoplastron of Ping-pien 18 lack numbers; the cracks for Ping-pien 20.7 and 20.8 lack numbers (cf. the k'ao-shih, p. 43), and so do at least five more inscriptionless cracks on the right and left hyo- and hypoplastrons of that shell. It is possible that in some of these cases the numbers may be on fragments now lost, but it is clear that numbers were not carved in every case (cf. ch. 2, n. 32, above). It appears that at the top center of the left hypoplastron of Ping-pien 14 the Shang engraver deliberately obliterated a crack number (k'ao-shih, p. 37); on this practice, see sec. 2.4. CCC C D U 0 81 }

\subsection[Interpretation]{\origsecnum{3.7.4}Interpretation}
\label{sec:chapters-ch03:3.7.4}

angle of the crack formed in hollow H15, which I have associated with subcharge 5: ``It
is because of Fu Hsin.'' The third crack notation is another hsiao-chi, ``slightly auspicious,''
just visible on the fragmented edge of the left hypoplastron of Ping-pien 18; this notation
lies in the lower angle of the crack in hollow H16, which I have associated with subcharge
6: It is not because of Fu Hsin.''
C C C D D
Interpretation
Having tentatively correlated the cracks, crack notations, and inscriptions, we
may now attempt to reconstruct the divination scenario. Two plastrons, the first and last
of the set, were sent in by Chih Fei at a place called Tuan; this fact was recorded as a
marginal notation. Other inscriptions reveal that the Shang farmed and sojourned at
Tuan and that it was the site of a garrison camp.<!-- FOOTNOTE:98 -->See Ping-pien 169.5-6; Yi-pien 1044; 1410 (S357.4). Since the two shells of our set were received at Tuan, it is conceivable that the divinations they record about the season's campaigns against the Hsia Wei and the Pa-fang were divined at ``camp Tuan,'' to which the king may have gone in the second month (Ping-pien 65.8 [S357.4]). The lack of any notation to the effect that our plastrons were ``ritually prepared'' might be explained by the fact that none of the royal women who usually performed this task (sec. 1.4) had accompanied the king to Tuan, where the shells were received.<!-- /FOOTNOTE --> .98 As a locus of such activities it was natural
that certain shipments of shell (sec. 1.2.4) would have been received there.<!-- FOOTNOTE:99 -->For other marginal notations recording the receipt of small numbers of shells at various places, see S256.<!-- /FOOTNOTE -->99 The source
of the other three shells in the set, one of which (Ping-pien 16/17) has been identified as
Ocadia sinensis (tables 3 and 4), was not recorded.<!-- FOOTNOTE:100 -->Another group, of three shells, was apparently received at Tuan (Yi-pien 642 [S357.4]), but the fragment recording that fact does not appear to belong to our set.<!-- /FOOTNOTE -->100
After the marginal notations had been carved in the shell (sec. 1.4), thirty-two double
hollows (sec. 1.5) were chiseled in the back of each of the five plastrons, which were then
ready for cracking. If we assume that it required thirty minutes to chisel a double hollow,
eighty man-hours would have been required to prepare the hollows on the five plastrons.<!-- FOOTNOTE:101 -->See the report of Chang Kuang-yüan, cited in ch. 1, n. 93.<!-- /FOOTNOTE --> 101
The Front of the Plastron
The first six charges concern contemplated attacks against two enemy statelets, the
Hsia Wei and the Pa-fang.<!-- FOOTNOTE:102 -->In the case of the Hsia Wei F, the reading of the second Shang graph, &, is uncertain; for an introduction to the scholarship and the location of the statelet, see Yü Hsing-wu (1940), pp. 22a-b; Shima (1958), pp. 389-390; Ikeda (1964), 1.18.3; Serruys (1974), p. 107, n. 39. In the case of the Pa-fang, there are disagreements about the reading of the first Shang graph, , and the location of the statelet; for an introduction to the problems, see Ch'en Meng-chia (1956), p. 284; Shima (1958), pp. 390-391; CKWT, p. 2783.<!-- /FOOTNOTE --> 102 The question in the king's mind was whether he should
follow Wang Cheng<!-- FOOTNOTE:103 -->Wang Ch'eng was a major Shang general or ally in period I. For an introduction to the scholarship, see Ogawa 1, shakubun, pp. 250-251. 83 CC C 0 U U<!-- /FOOTNOTE -->103 in his attack on the Hsia Wei or whether he should follow Chih

\subsubsection[The Front of the Plastron]{\origsecnum{3.7.4.1}The Front of the Plastron}
\label{sec:chapters-ch03:3.7.4.1}

angle of the crack formed in hollow H15, which I have associated with subcharge 5: ``It
is because of Fu Hsin.'' The third crack notation is another hsiao-chi, ``slightly auspicious,''
just visible on the fragmented edge of the left hypoplastron of Ping-pien 18; this notation
lies in the lower angle of the crack in hollow H16, which I have associated with subcharge
6: It is not because of Fu Hsin.''
C C C D D
Interpretation
Having tentatively correlated the cracks, crack notations, and inscriptions, we
may now attempt to reconstruct the divination scenario. Two plastrons, the first and last
of the set, were sent in by Chih Fei at a place called Tuan; this fact was recorded as a
marginal notation. Other inscriptions reveal that the Shang farmed and sojourned at
Tuan and that it was the site of a garrison camp.<!-- FOOTNOTE:98 -->See Ping-pien 169.5-6; Yi-pien 1044; 1410 (S357.4). Since the two shells of our set were received at Tuan, it is conceivable that the divinations they record about the season's campaigns against the Hsia Wei and the Pa-fang were divined at ``camp Tuan,'' to which the king may have gone in the second month (Ping-pien 65.8 [S357.4]). The lack of any notation to the effect that our plastrons were ``ritually prepared'' might be explained by the fact that none of the royal women who usually performed this task (sec. 1.4) had accompanied the king to Tuan, where the shells were received.<!-- /FOOTNOTE --> .98 As a locus of such activities it was natural
that certain shipments of shell (sec. 1.2.4) would have been received there.<!-- FOOTNOTE:99 -->For other marginal notations recording the receipt of small numbers of shells at various places, see S256.<!-- /FOOTNOTE -->99 The source
of the other three shells in the set, one of which (Ping-pien 16/17) has been identified as
Ocadia sinensis (tables 3 and 4), was not recorded.<!-- FOOTNOTE:100 -->Another group, of three shells, was apparently received at Tuan (Yi-pien 642 [S357.4]), but the fragment recording that fact does not appear to belong to our set.<!-- /FOOTNOTE -->100
After the marginal notations had been carved in the shell (sec. 1.4), thirty-two double
hollows (sec. 1.5) were chiseled in the back of each of the five plastrons, which were then
ready for cracking. If we assume that it required thirty minutes to chisel a double hollow,
eighty man-hours would have been required to prepare the hollows on the five plastrons.<!-- FOOTNOTE:101 -->See the report of Chang Kuang-yüan, cited in ch. 1, n. 93.<!-- /FOOTNOTE --> 101
The Front of the Plastron
The first six charges concern contemplated attacks against two enemy statelets, the
Hsia Wei and the Pa-fang.<!-- FOOTNOTE:102 -->In the case of the Hsia Wei F, the reading of the second Shang graph, &, is uncertain; for an introduction to the scholarship and the location of the statelet, see Yü Hsing-wu (1940), pp. 22a-b; Shima (1958), pp. 389-390; Ikeda (1964), 1.18.3; Serruys (1974), p. 107, n. 39. In the case of the Pa-fang, there are disagreements about the reading of the first Shang graph, , and the location of the statelet; for an introduction to the problems, see Ch'en Meng-chia (1956), p. 284; Shima (1958), pp. 390-391; CKWT, p. 2783.<!-- /FOOTNOTE --> 102 The question in the king's mind was whether he should
follow Wang Cheng<!-- FOOTNOTE:103 -->Wang Ch'eng was a major Shang general or ally in period I. For an introduction to the scholarship, see Ogawa 1, shakubun, pp. 250-251. 83 CC C 0 U U<!-- /FOOTNOTE -->103 in his attack on the Hsia Wei or whether he should follow Chih

\section[Conclusions]{\origsecnum{3.8}Conclusions}
\label{sec:chapters-ch03:3.8}

I do not insist upon the accuracy of the above interpretation in every detail. The
meaning of some of the charges and terms used is still obscure; and the relation between
the charges and the cracks remains hypothetical. But this kind of attempt to read an
inscription in full context and to discover its rationale must be undertaken, whenever
possible, before an inscription is used as a historical source. Such close analysis, with one's
nose to the bone as it were, may appear tedious at times-it is undoubtedly tedious to
read about it in this way—but it is the kind of analysis that will, on occasion, provide
insights that advance our knowledge of Shang divination and culture in fresh and
unexpected ways. Scholars must understand the other inscriptions on the same bone
or shell, both front and back, and the possible relations between them; they must look
for crack notations; they must, using crack numbers, look for other inscriptions from the
same set; and they must, using Shima's Sōrui, look for related inscriptions. They must try
to discover what was going on-on the oracle bone, in the Shang court, in the diviner's
mind.

